\SourcePage{126}
\SGAchapter{XIV}{与超越理论的比较}{Comparaison avec la théorie transcendante}
\begin{center}P. Deligne\end{center}
\textbf{引言。}
为了把第 XIII 讲的\Term{vanishing-cycles}{消失圈}理论与其在 $\mathbf C$ 上的\Term{transcendental}{超越}对应理论相比较，我们分两步进行。

 a) 我们将仿照第 XIII 讲的\Term{form}{形式}体系，给出一个\Term{transcendental}{超越}\Term{form}{形式}体系，并用通常的方法证明一个\Term{comparison-theorem}{比较定理}（这些方法已用于证明第 XIII 讲 2.3 的\Term{finiteness-theorem}{有限性定理}）。

 b) 在\Term{isolated-singularity}{孤立奇点}的情形，我们将用更具体的语言阐明所引入的\Term{transcendental}{超越}\Term{form}{形式}体系的含义。

除\Term{comparison-theorem}{比较定理}外，本讲还给出 \Term{picard-lefschetz-formula}{Picard--Lefschetz 公式}在\Term{transcendental}{超越}情形的一个证明。
\SourcePage{127}
在第 4 节，我们给出\Term{vanishing-cycles}{消失圈}的 \Term{hodge-theory}{Hodge 理论}的一个预备结果。本讨论班后续部分不使用该节。

\SGAsection{1}{消失圈的超越形式体系}
\textbf{1.1. \Term{punctured-disc}{穿孔圆盘}的\Term{fundamental-group}{基本群}：\Term{signature}{符号}约定。}

我们将要考虑\Term{punctured-disc}{穿孔圆盘} $D^*$ 及其\Term{fundamental-group}{基本群} $\mathbf Z$。这个群是交换的，因而可能出现\Term{signature}{符号}歧义。忘掉它的\Term{commutativity}{交换性}，就可以消除这些歧义。

\textbf{1.1.1.} 对于良好的\Term{connected}{连通}\Term{topological-space}{拓扑空间} $X$ 及 $x\in X$，\Term{fundamental-group}{基本群}的元素是以 $x$ 为起点的\Term{loop}{闭路}的\Term{homotopy}{同伦}类；若 $\alpha,\beta\in\pi_1(X,x)$ 分别由\Term{loop}{闭路} $a,b$ 表示，则 $\alpha\beta$ 由先走 $b$ 再走 $a$ 所得到的\Term{loop}{闭路}表示。

\textbf{1.1.2.} 设 $F$ 是 $X$ 上的\Term{locally-constant-sheaf}{局部常值层}。对每条\Term{path}{道路} $a:[0,1]\to X$，$F$ 在 $[0,1]$ 上的\Term{inverse-image}{逆像} $a^*F$ 是\Term{locally-constant-sheaf}{局部常值层}，从而是\Term{constant-sheaf}{常值层}，故有\Term{bijection}{双射}
\[
 a:F_{a(0)}\xleftarrow{\sim}H^0([0,1],a^*F)\xrightarrow{\sim}F_{a(1)},
\]
它只依赖于 $a$ 的\Term{homotopy}{同伦}类。这一构造给出 $\pi_1(X,x)$ 在 $F_x$ 上的（左）\Term{action}{作用}。\Term{functor}{函子} $F\mapsto F_x$ 是 $X$ 上\Term{locally-constant-sheaf}{局部常值层}的\Term{category}{范畴}与带有 $\pi_1(X,x)$ \Term{action}{作用}的\Term{set}{集合}的\Term{category}{范畴}之间的一个\Term{equivalence}{等价}。

\textbf{1.1.3.} 当 $x$ 变化时，各 $\pi_1(X,x)$ 是一个\Term{locally-constant-sheaf}{局部常值群层} $\underline\pi_1(X)$ 的\Term{fibre}{纤维}；这个\Term{sheaf}{层}\Term{action}{作用}于每个\Term{locally-constant-sheaf}{局部常值层} $F$。若 $\pi_1(X,x)$ 是\Term{abelian-group}{交换群}，则 $\underline\pi_1(X)$ 是\Term{constant-sheaf}{常值层}，其值记为 $\pi_1(X)$。

\textbf{1.1.4.} 群 $\pi_1(X,x)$ \Term{action}{作用}于 $(X,x)$ 的\Term{universal-covering}{万有覆盖} $p:(\widetilde X,\widetilde x)\to(X,x)$：若把这个\Term{covering}{覆盖}看作 $X$ 上的\Term{etale-space}{平展空间}，从而看作一个\Term{sheaf}{层}（其\Term{local-scheme}{局部}\Term{section}{截面}的\Term{sheaf}{层}），则由 1.1.2，它对应于带有\Term{identity-element}{单位元} $e$ 的 $\pi_1(X,x)$，
\SourcePage{128}
其上 $\pi_1(X,x)$ 按\Term{left-translation}{左平移}\Term{action}{作用}；群 $\pi_1(X,x)$ \Term{action}{作用}于这个 $\pi_1$ \Term{set}{集合}（通过 $\pi_1$ \Term{set}{集合}的\Term{automorphism}{自同构}）：$\sigma\in\pi_1(X,x)$ 的\Term{action}{作用}为 $x\mapsto x\sigma^{-1}$。

这个构造把 $\pi_1(X,x)$ 与 $\operatorname{Aut}_X(\widetilde X)$ 等同起来。

若 $F$ 是 $X$ 上的\Term{locally-constant-sheaf}{局部常值层}，则\Term{bijection}{双射}
\[
 H^0(\widetilde X,p^*F)\longrightarrow(p^*F)_{\widetilde x}=F_x
 \tag{1.1.4.1}
\]
与 $\pi_1(X,x)$ 的\Term{action}{作用}交换：左端的\Term{action}{作用}通过它在 $\widetilde X$ 上的\Term{action}{作用}作结构搬移得到，右端的\Term{action}{作用}由 1.1.2 给出。1.1.2 的\Term{equivalence}{等价}也可以表述为：
\[
 F\longmapsto\bigl(\widetilde X\longmapsto H^0(\widetilde X,p^*F)\bigr)
\]
是 $X$ 上\Term{locally-constant-sheaf}{局部常值层}的\Term{category}{范畴}与\Term{functor}{函子}
\[
 (X\text{ 的万有覆盖})\longrightarrow(\mathrm{Ens})
\]
的\Term{category}{范畴}之间的一个\Term{equivalence}{等价}。

\textbf{1.1.5.} 若 $\pi_1(X,x)$ 是\Term{abelian-group}{交换群}，因而与 $x$ 无关，则 1.1.4 中 $\pi_1(X)$ 在 $\widetilde X$ 上的\Term{action}{作用}与\Term{base-point}{基点}无关：$\pi_1(X)$ \Term{action}{作用}于 $X$ 的每个\Term{universal-covering}{万有覆盖} $\widetilde X$。

要注意，这个\Term{action}{作用}（以后只考虑此\Term{action}{作用}）对应于 1.1.3 中\Term{action}{作用}的逆。

\textbf{1.1.6.} 设 $D$ 是一个\Term{disc}{圆盘}，即与标准\Term{disc}{圆盘} $\{z\in\mathbf C\mid |z|<1\}$ \Term{homeomorphism}{同胚}的\Term{topological-space}{拓扑空间}。设 $0\in D$，$D^*=D-\{0\}$。假定 $D$ 带有\Term{analytic-manifold}{解析流形}结构，因而具有\Term{orientation}{定向}。于是\Term{fundamental-group}{基本群} $\pi_1(D^*)$ \Term{canonical}{典范}地\Term{isomorphism}{同构}于 $\mathbf Z$：对于标准\Term{disc}{圆盘}、$0=0$ 及 $t\in D^*$，取\Term{loop}{闭路} $x\mapsto\exp(2\pi i x)t$ 为 $\pi_1(D^*,t)$ 的正\Term{generator}{生成元}。

\textbf{1.1.7.} 设 $z:D\to\mathbf C$ 是一个\Term{uniformizer}{局部参数}（$(dz)_0\ne0$）。设 $\widetilde D^*$ 是 $D^*$ 的一个\Term{universal-covering}{万有覆盖}。用 $\log z$ 表示 $\widetilde D^*$ 上任一满足
\SourcePage{129}
$\exp(\log z)=z$ 的函数。二元组 $(\widetilde D^*,\log z)$ 在唯一\Term{isomorphism}{同构}意义下唯一。

对 $\alpha\in\mathbf C$，置 $z^\alpha=\exp(\alpha\log z)$。若 $T$ 为 $\pi_1(D^*)$ 的正\Term{generator}{生成元}，它由 1.1.5 \Term{action}{作用}于 $\widetilde D^*$，并通过结构搬移\Term{action}{作用}于 $\widetilde D^*$ 上的函数，则有
\[
 T\log(z)=\log(z)+2\pi i,\tag{1.1.7.1}
\]
\[
 Tz^\alpha=\exp(2\pi i\alpha)z^\alpha.\tag{1.1.7.2}
\]

\textbf{1.1.8.} 设 $D$ 是一个\Term{disc}{圆盘}，$0\in D$，$D^*=D-\{0\}$。若 $D'\subset D$ 是包含 $0$ 的较小\Term{disc}{圆盘}，则\Term{restriction}{限制}到 $D'^*$ 给出 $D^*$ 上\Term{locally-constant-sheaf}{局部常值层}的\Term{category}{范畴}与 $D'^*$ 上\Term{locally-constant-sheaf}{局部常值层}的\Term{category}{范畴}之间的\Term{equivalence}{等价}。事实上，对 $t\in D'^*$，有 $\pi_1(D'^*,t)\xrightarrow{\sim}\pi_1(D^*,t)$。

假定 $D$ 带有\Term{differentiable-manifold}{微分流形}结构。设 $T$ 是 $D$ 在 $0$ 处的\Term{tangent-space}{切空间}，$T^*=T-\{0\}$。对上述 $D'\subset D$ 及满足 $f(0)=0$、$(df)_0=\mathrm{Id}$ 的\Term{open-embedding}{开嵌入} $f:D'\hookrightarrow T$，复合
\[
 (\text{限制到 }D')^{-1}\circ f^*
\]
是 $T^*$ 上\Term{locally-constant-sheaf}{局部常值层}的\Term{category}{范畴}与 $D^*$ 上\Term{locally-constant-sheaf}{局部常值层}的\Term{category}{范畴}之间的一个\Term{equivalence}{等价} $[f]$。在\Term{canonical}{典范}\Term{isomorphism}{同构}意义下，它只依赖于 $f$ 在 $0$ 处的\Term{germ}{芽}，该\Term{germ}{芽}满足 $(df)_0=\mathrm{Id}$。从一个这样的\Term{germ}{芽}到另一个的\Term{path}{道路}定义一个\Term{equivalence}{等价}之间的\Term{isomorphism}{同构}，而两条\Term{homotopy}{同伦}的\Term{path}{道路}定义相同的\Term{isomorphism}{同构}。所考虑的\Term{germ}{芽}的\Term{set}{集合}是\Term{contractible}{可缩}的（向 $f_0$ 的\Term{retraction}{收缩}为 $(1-\lambda)f+\lambda f_0$），故得到 $T^*$ 上\Term{locally-constant-sheaf}{局部常值层}的\Term{category}{范畴}与 $D^*$ 上\Term{locally-constant-sheaf}{局部常值层}的\Term{category}{范畴}之间的一个\Term{equivalence}{等价}，它在唯一\Term{isomorphism}{同构}意义下确定。

\textbf{1.1.9.} 对一个\Term{disc}{圆盘} $D$、$0\in D$、$D^*=D-\{0\}$ 及 $D^*$ 的\Term{universal-covering}{万有覆盖} $p:\widetilde D^*\to D^*$，我们用 $p:\widetilde D\to D$ 表示在 $\widetilde D^*$ 中添入一个点 $0$ 所得的 $D$ 上的\Term{topological-space}{拓扑空间}，使得
\SourcePage{130}

 a) $\widetilde D^*$ 是 $\widetilde D$ 中的\Term{open-subset}{开集}；

 b) $p:\widetilde D\to D$ 延拓 $p:\widetilde D^*\to D^*$，且 $p(0)=0$；

 c) 当 $U$ 遍历 $D$ 中 $0$ 的\Term{neighbourhood}{邻域}时，$p^{-1}(U)$ 构成 $\widetilde D$ 中 $0$ 的一个\Term{neighbourhood-basis}{基本邻域系}。

\textbf{1.1.10.} 设 $X$ 是一条\Term{nonsingular}{非奇异}复\Term{algebraic-curve}{代数曲线}，$0\in X$，$D$ 是 $X$ 中 $0$ 的一个\Term{disc}{圆盘}\Term{neighbourhood}{邻域}。$D$ 上在 $X$ 的\Term{rational-function-field}{有理函数域} $K(X)$ 上\Term{algebra}{代数}的\Term{holomorphic}{全纯}函数在 $0$ 处的\Term{germ}{芽}所成的环，是 $X$ 在 $0$ 处的\Term{local-ring}{局部环} $\mathscr O_0$ 的\Term{henselization}{Hensel 化} $\mathscr O_0^h$。设 $\widetilde D^*$ 是 $D^*=D-\{0\}$ 的一个\Term{universal-covering}{万有覆盖}。$\widetilde D^*$ 上在 $K(X)$ 上\Term{algebra}{代数}的\Term{meromorphic}{亚纯}函数在 $0\in\widetilde D$ 处的\Term{germ}{芽}所成的域，是 $K(X)$ 的一个\Term{algebraic-closure}{代数闭包} $\overline K$（从而也是 $\mathscr O_0^h$ 的\Term{fraction-field}{分式域} $K_0^h$ 的\Term{algebraic-closure}{代数闭包}）。

由 $\widetilde D^*$ 的 $D$-\Term{automorphism}{自同构}组成的群 $\pi_1(D^*)$（1.1.5）通过结构搬移\Term{action}{作用}于 $\overline K$，因而有
\[
\begin{tikzcd}
\pi_1(D^*)\arrow[r]\arrow[d,equal]&\operatorname{Gal}(\overline K/K_0^h)\arrow[d,equal]\\
\mathbf Z&\widehat{\mathbf Z}(1).
\end{tikzcd}\tag{1.1.10.1}
\]
由 (1.1.7.2)，此图与 $\mathbf Z$ 到 $\widehat{\mathbf Z}(1)$ 的自然\Term{embedding}{嵌入}（在 $\mathbf C$ 上）相容；后者是\Term{embedding}{嵌入} $\mathbf Z\hookrightarrow\widehat{\mathbf Z}$ 与\Term{isomorphism}{同构} $\widehat{\mathbf Z}\xrightarrow{\sim}\widehat{\mathbf Z}(1)$ 的复合，此\Term{isomorphism}{同构}是各\Term{isomorphism}{同构}
\[
 k\longmapsto\exp\left(\frac{2\pi i}{n}k\right):\mathbf Z/n\longrightarrow\mu_n
\]
的\Term{inverse-limit}{投射极限}。

\textbf{1.1.11.} $\operatorname{Spec}(K_0^h)$ 上的每个\Term{sheaf}{层} $F$ 都定义 $D^*$ 上的一个\Term{locally-constant-sheaf}{局部常值层}：利用 1.1.4 的\Term{equivalence}{等价}，给 $D^*$ 的每个\Term{universal-covering}{万有覆盖} $\widetilde D^*$ 配上 $F_{\overline K}$，其中 $\overline K$ 由 1.1.10 构造。

\textbf{1.2. \Term{set-sheaf}{集合层}。}

\textbf{1.2.1.} 设 $D$ 是一个\Term{disc}{圆盘}（1.1.6），$0\in D$，$D^*=D-\{0\}$。我们用 $\mathscr D^*$ 表示
\SourcePage{131}
$D^*$ 上\Term{locally-constant-sheaf}{局部常值集合层}的\Term{topos}{拓扑斯}。由 1.1.8，若把 $D$ 换成包含 $0$ 的较小\Term{disc}{圆盘} $D'$，或者当 $D$ 带有\Term{differentiable-manifold}{微分流形}结构时把 $D$ 换成它在 $0$ 处的\Term{tangent-space}{切空间}，这个\Term{topos}{拓扑斯}都不“改变”。

对 $t\in D^*$，\Term{functor}{函子} $F\mapsto F_t$ 把这个\Term{topos}{拓扑斯}与带有 $\pi_1(D^*)$ \Term{action}{作用}的\Term{set}{集合}的\Term{topos}{拓扑斯}等同起来（1.1.2）。由 1.1.4，\Term{functor}{函子}
\[
 F\longmapsto\bigl(\widetilde D^*\longmapsto H^0(\widetilde D^*,p^*F)\bigr)
\]
把它与\Term{functor}{函子}
\[
 (D^*\text{ 的万有覆盖})\longrightarrow(\mathrm{Ens})
\]
的\Term{topos}{拓扑斯}等同起来，而这些等同之间的关系由 (1.1.4.1) 给出。

\textbf{1.2.2.} 设 $F$ 是 $D$ 上的一个\Term{sheaf}{层}，它在 $D^*$ 上的\Term{restriction}{限制}是\Term{locally-constant-sheaf}{局部常值层}。存在唯一的\Term{sheaf}{层}\Term{morphism}{态射} $F_0\xrightarrow{\alpha}F$，在 $0$ 处为恒等。如第 XIII 讲 0.1 中所说明的（取 $X_0=(\mathrm{pt})$），\Term{functor}{函子}
\[
 F\longmapsto(F_0,F\text{ 在 }D^*\text{ 上的限制},\alpha)
\]
把这些\Term{sheaf}{层}等同于如下三元组：一个\Term{set}{集合} $F_0$、$D^*$ 上的一个\Term{locally-constant-sheaf}{局部常值层} $F_\eta$，以及
\[
 \alpha:(D^*\text{ 上值为 }F_0\text{ 的常值层})\longrightarrow(\text{层 }F_\eta).
\]

\textbf{1.2.3.} 用 $\mathscr D$ 表示 $D$ 上\Term{restriction}{限制}到 $D^*$ 为\Term{locally-constant-sheaf}{局部常值层}的\Term{set}{集合}\Term{sheaf}{层}的\Term{topos}{拓扑斯}。$\mathscr D$ 到 $D$ 上\Term{sheaf}{层}中的包含，是一个\Term{topos}{拓扑斯}\Term{morphism}{态射} $D\to\mathscr D$ 的\Term{inverse-image}{逆像}\Term{functor}{函子}；它\Term{embedding}{嵌入}如下\Term{commutative-diagram}{交换图}：
\[
\begin{tikzcd}
\{0\}\arrow[r,hook]\arrow[d,equal]&D\arrow[d,"\beta"]&D^*\arrow[l,hook']\arrow[d,"\beta"]\\
\{0\}\arrow[r,hook]&\mathscr D&\mathscr D^*\arrow[l,hook'].
\end{tikzcd}
\]
如同 1.1.1，可以验证，把 $D$ 换成较小\Term{disc}{圆盘}或换成 $D$ 在 $0$ 处的\Term{tangent-space}{切空间}，此\Term{topos}{拓扑斯}\Term{invariant}{不变}。

\SourcePage{132}
\textbf{1.2.4.} 设 $Y$ 是一个\Term{topological-space}{拓扑空间}。\Term{product-topos}{乘积拓扑斯} $Y\times\mathscr D$ 有下列三种描述（其\Term{equivalence}{等价}性见第 XIII 讲 0.1）。

 a) $Y\times\mathscr D$ 是满足下列条件的 $Y\times D$ 上的\Term{sheaf}{层} $F$ 的\Term{topos}{拓扑斯}：

 $(*)$ $D^*$ 可由一些\Term{open-subset}{开集} $U$ \Term{covering}{覆盖}，使 $F|_{Y\times U}$ 是 $Y$ 上某个\Term{sheaf}{层}的\Term{inverse-image}{逆像}。

 b) 设 $t\in D^*$。$Y\times\mathscr D$ 是三元组 $(F_0,F_t,\alpha)$ 的\Term{topos}{拓扑斯}，其中：

 $\alpha)$ $F_0$ 是 $Y$ 上的\Term{sheaf}{层}；

 $\beta)$ $F_t$ 是 $Y$ 上带有 $\pi_1(D^*)$ \Term{action}{作用}的\Term{sheaf}{层}；

 $\gamma)$ $\alpha$ 是从 $F_0$ 到 $F_t$ 的\Term{morphism}{态射}，其像包含于\Term{invariant}{不变}元中。

 c) $Y\times\mathscr D$ 是三元组 $(F_0,F_\eta,\alpha)$ 的\Term{topos}{拓扑斯}，其中：

 $\alpha')$ $F_0$ 是 $Y$ 上的\Term{sheaf}{层}；

 $\beta')$ $F_\eta$ 是如下\Term{functor}{函子}，其值记为 $F_{\widetilde\eta}$：
\[
 (D^*\text{ 的万有覆盖})\longrightarrow(Y\text{ 上的层});
\]
 $\gamma')$ $\alpha$ 是从 $F_0$（从 $D^*$ 的\Term{universal-covering}{万有覆盖}到 $Y$ 上\Term{sheaf}{层}的、值为 $F_0$ 的常\Term{functor}{函子}）到 $F_\eta$ 的\Term{morphism}{态射}。

此外，这个\Term{topos}{拓扑斯}是否确为 $Y$ 与 $\mathscr D$ 的\Term{product-topos}{乘积拓扑斯}，对我们并不重要。类似地，可以任选下述描述来定义 $Y\times\mathscr D^*$：

 a$^*$) 满足 $(*)$ 的 $Y\times D^*$ 上\Term{sheaf}{层}的\Term{topos}{拓扑斯}；

 b$^*$)（对 $t\in D^*$）$Y$ 上带有 $\pi_1(D^*)$ \Term{action}{作用}的\Term{sheaf}{层} $F_t$ 的\Term{topos}{拓扑斯}；

 c$^*$) \Term{functor}{函子}
\[
 (D^*\text{ 的万有覆盖})\longrightarrow(Y\text{ 上的层})
\]
的\Term{topos}{拓扑斯}。

\textbf{1.2.5.} \Term{topos}{拓扑斯} $Y\times\mathscr D$ 关于 $Y$ 的\Term{functoriality}{函子性}，与第 XIII 讲的\Term{topos}{拓扑斯} $Y\times_s S$ 相同，我们不再逐一列举。为固定记号，只给出图
\SourcePage{133}
\[
\begin{tikzcd}
Y\arrow[d]&Y\arrow[l,"\mathrm{sp}"]\arrow[r,hook,"i"]\arrow[d]&Y\times\mathscr D\arrow[d]&Y\times\mathscr D^*\arrow[l,hook',"j"']\arrow[d]\\
(\mathrm{pt})&\{0\}\arrow[r,hook,"i"]&\mathscr D&\mathscr D^*\arrow[l,hook',"j"'].
\end{tikzcd}
\]

\textbf{1.2.6.} 设 $K$ 是 $Y\times\mathscr D$ 上的一个阿贝尔\Term{sheaf}{层}\Term{complex}{复形}。仿照第 XIII 讲 1.4.2，从 $K$ 得到 $\mathrm K(Y\times\mathscr D^*)$ 中的一个\Term{distinguished-triangle}{区别三角}
\[
\begin{tikzcd}[column sep=large]
&\Phi(K)\arrow[dl]&\\
\mathrm{sp}^*K_0\arrow[rr]&&K_\eta\arrow[ul].
\end{tikzcd}
\]
这个构造\Term{descent}{下降}到\Term{derived-category}{导出范畴}，并定义一个\Term{functor}{函子}
\[
 \mathrm D(Y\times\mathscr D)\longrightarrow
 (\mathrm D(Y\times\mathscr D^*)\text{ 中的区别三角}).
\]
假定 $D$ 已\Term{orientation}{定向}（例如是复解析的），设 $T$ 为 $\pi_1(D^*)$ 的正\Term{generator}{生成元}（1.1.6）。设 $K\in\operatorname{Ob}\mathrm D(Y\times\mathscr D)$。仿照第 XIII 讲 1.4.3，定义\Term{variation}{变差} $V=\operatorname{Var}(T)$：
\[
 V:\Phi(K)\longrightarrow K_\eta.
\]
$K_\eta$ 的\Term{endomorphism}{自态射} $T-\mathrm{Id}$ 是复合
\[
 T-\mathrm{Id}:K_\eta\longrightarrow\Phi(K)\longrightarrow K_\eta.
\]

\textbf{1.2.7.} 设 $D$ 为\Term{disc}{圆盘}，$0\in D$，$D^*=D-\{0\}$，$f:X\to D$ 为\Term{continuous-map}{连续映射}，$X_0=f^{-1}(0)$。我们要定义一个\Term{functor}{函子}
\[
 \Psi:(X\text{ 上的层})\longrightarrow(X_0\times\mathscr D\text{ 上的层}).
\]
设 $\widetilde D^*$ 为 $D^*$ 的一个\Term{universal-covering}{万有覆盖}，$\widetilde D$ 如 1.1.9 所述，
\SourcePage{134}
\[
 \overline X=X\times_D\widetilde D,\qquad
 X^*=X-X_0=X\times_D D^*,\qquad
 \widetilde X^*=\overline X-X_0=X\times_D\widetilde D^*.
\]
有\Term{commutative-diagram}{交换图}
\[
\begin{tikzcd}
X_0\arrow[r,hook,"\overline i"]\arrow[d,equal]&\overline X\arrow[d,"p"]&\widetilde X^*\arrow[l,hook',"\overline j"']\arrow[d,"p'"]\\
X_0\arrow[r,hook,"i"]&X&X^*\arrow[l,hook',"j"']
\end{tikzcd}\tag{1.2.7.1}
\]
它位于图
\[
\begin{tikzcd}
\{0\}\arrow[r,hook,"\overline i"]\arrow[d,equal]&\widetilde D\arrow[d,"p"]&\widetilde D^*\arrow[l,hook',"\overline j"']\arrow[d,"p'"]\\
\{0\}\arrow[r,hook,"i"]&D&D^*\arrow[l,hook',"j"']
\end{tikzcd}\tag{1.2.7.2}
\]
之上，并且是\Term{cartesian-diagram}{笛卡尔}的。

对 $X$ 上的\Term{sheaf}{层} $F$，置（使用 1.2.4 c) 的记号）
\[
 \Psi(F)=\bigl(i^*F,\overline i^*\overline j_*((jp')^*F),\alpha\bigr),
\]
其中
\[
 \alpha:i^*F=\overline i^*p^*F\longrightarrow
 \overline i^*\overline j_*((jp')^*F)
 =\overline i^*\overline j_*\overline j^*p^*F
\]
由\Term{adjunction}{伴随}\Term{morphism}{态射} $\mathrm{Id}\to\overline j_*\overline j^*$ 经 $\overline i^*$ 得到。

还定义\Term{functor}{函子}
\[
 \Psi_\eta:(X^*\text{ 上的层})\longrightarrow
 (X_0\times\mathscr D^*\text{ 上的层}),\qquad
 \Psi_\eta(F)=\overline i^*\overline j_*(p'^*F).
\]

\textbf{1.2.8.} $\Psi$ 的\Term{functoriality}{函子性}，与第 XIII 讲 1.3.5 至 1.3.9 中研究的对应\Term{functor}{函子}相同。

\SourcePage{135}
\textbf{1.3. \Term{functor}{函子} $R\Psi$。}

\textbf{1.3.1.} 设 $D$ 是\Term{disc}{圆盘}，$0\in D$，$D^*=D-\{0\}$，$f:X\to D$ 是\Term{continuous-map}{连续映射}，$X_0=f^{-1}(0)$。\Term{functor}{函子} $R\Psi$ 与 $R\Psi_\eta$ 是\Term{functor}{函子} $\Psi$ 与 $\Psi_\eta$ 的\Term{derived-functor}{导出函子}：
\[
\begin{aligned}
 R\Psi&:\mathrm D^+(X)\longrightarrow\mathrm D^+(X_0\times\mathscr D),\\
 R\Psi_\eta&:\mathrm D^+(X^*)\longrightarrow\mathrm D^+(X_0\times\mathscr D^*).
\end{aligned}
\]
以 $R\Phi$ 表示复合\Term{functor}{函子} $\Phi\circ R\Psi$：
\[
 R\Phi:\mathrm D^+(X)\longrightarrow\mathrm D^+(X_0\times\mathscr D^*).
\]
这些\Term{functor}{函子}的性质与第 XIII 讲 2.1 中的同名\Term{functor}{函子}完全类似。其证明初等得多，因为在\Term{topology}{拓扑}框架中，\Term{proper-morphism}{固有态射}或\Term{smooth-morphism}{光滑态射}（即\Term{local-scheme}{局部}\Term{isomorphism}{同构}于 $\mathbf R^n\times S\to S$）的\Term{base-change}{换基}\Term{theorem}{定理}都很容易证明。下面列出所需作的少量改动。

\textbf{1.3.2.} \Term{functor}{函子} $Rf_!$ 只对\Term{locally-compact}{局部紧}\Term{topological-space}{空间}之间的\Term{separated-morphism}{分离态射}定义，它是\Term{proper-support}{固有支撑}\Term{direct-image}{正像}\Term{functor}{函子} $f_!$ 的\Term{derived-functor}{导出函子}。

\textbf{1.3.3.} 对\Term{proper-morphism}{固有}（且\Term{separated-morphism}{分离}）\Term{morphism}{态射} $f:X\to D$，基本\Term{isomorphism}{同构} (2.1.8.3) 应改写如下。设 $p:\widetilde D^*\to D^*$ 是 $D^*$ 的\Term{universal-covering}{万有覆盖}。对 $D$ 中 $0$ 的每个\Term{neighbourhood}{邻域} $U$，设 $\widetilde U^*=p^{-1}(U)$，$X_{\widetilde U^*}=X\times_D\widetilde U^*$。对 $K\in\operatorname{Ob}\mathrm D^+(f^{-1}(D^*))$，有
\[
 \varinjlim_U\mathbb H^i(X_{\widetilde U^*},K)
 \xrightarrow{\sim}\mathbb H^i(X_0,R\Psi_\eta(K)).\tag{1.3.3.1}
\]
这个\Term{isomorphism}{同构}在 $X^*=f^{-1}(D^*)$ 是 $D^*$ 上的\Term{topological-fibration}{拓扑纤维丛}时尤其有用。当 $U\subset D$ 为趋向 $0$ 的\Term{disc}{圆盘}时，
\SourcePage{136}
$X_{\widetilde U^*}\to\widetilde U^*$ 的 Leray \Term{spectral-sequence}{谱序列}表明，对 $X^*$ 上的每个\Term{locally-constant-sheaf}{局部常值层} $\Lambda$ 及每个 $t\in\widetilde U^*$，有
\[
 H^i(X_{\widetilde U^*},\Lambda)\xrightarrow{\sim}H^i(X_t,\Lambda).
\]
因此，对 $t\in D^*$，有
\[
 H^i(X_t,\Lambda)\xrightarrow{\sim}H^i(X_0,R\Psi_t(\Lambda))\tag{1.3.3.2}
\]
（$X/D$ \Term{proper-morphism}{固有}，$X^*$ 是\Term{topological-fibration}{纤维丛}，$\Lambda$ \Term{locally-constant-sheaf}{局部常值}）。

实际中，为验证 $X^*$ 是 $D^*$ 上的\Term{topological-fibration}{纤维丛}，使用 Thom 的\Term{isotopy}{同痕}\Term{theorem}{定理}，下面给出其一种\Term{form}{形式}。（此\Term{theorem}{定理}的证明参见 Whitney 和 Mather；它的任何推论都不会在以后的各讲中使用。）

\textbf{1.3.4. 预备。} 设 $f:X\to Y$ 是\Term{differentiable-manifold}{微分流形}的\Term{smooth-morphism}{光滑态射}（按定义，即没有\Term{critical-point}{临界点}的\Term{morphism}{态射}）。$X/Y$ 上的一个\Term{connection}{联络} $\nabla$，是 $X$ 的\Term{tangent-bundle}{切丛}中与 $\operatorname{Ker}(df)$ 互补的一个\Term{vector-subbundle}{向量子丛}。对 $Y$ 上的\Term{vector-field}{向量场} $u$，在此\Term{vector-subbundle}{子丛}中存在唯一一个\Term{lifting}{提升} $u$ 的 $X$ 上\Term{vector-field}{向量场} $\nabla(u)$。若
\[
 \nabla([u,v])=[\nabla(u),\nabla(v)],
\]
则称 $\nabla$ 是\Term{integrable-connection}{可积联络}。当 $f$ \Term{proper-morphism}{固有}时，一个\Term{integrable-connection}{可积联络}在 $Y$ 上定义 $X/Y$ 的\Term{local-scheme}{局部}\Term{trivialization}{平凡化}
\[
 \alpha:f^{-1}(U)\xrightarrow{\sim}X_0\times U,
\]
它与到 $U$ 的\Term{projection}{投影}相容，且 $d\alpha(\nabla(u))=(0,u)$。当 $f$ 不\Term{proper-morphism}{固有}时，在无穷远处有适当条件下仍有同样结果，例如存在适当的\Term{exhaustion-function}{穷竭函数} $g$，使各 $\nabla(u)g$ 为零。

\begin{SGAresult}{Thom \Term{isotopy}{同痕}\Term{theorem}{定理}}{1.3.5}
设 $D=\{z\in\mathbf C\mid |z|<1\}$，$D^*=D-\{0\}$，$f:X\to D$ 是\Term{analytic-space}{解析空间}的一个\Term{proper-morphism}{固有态射}。把 $D$ 换成较小的\Term{disc}{圆盘}，并把 $X$ 换成它的\Term{inverse-image}{逆像}后，存在 $X^*=f^{-1}(D^*)$ 的一个\Term{stratification}{分层} $\Sigma$，具有下列性质。

(a) $\Sigma$ 是 $X^*$ 的\Term{closed-subset}{闭子集}\Term{filtration}{滤过}
\[
 X^*=F_0\supset F_1\supset\cdots\supset F_n\supset F_{n+1}=\varnothing,
\]
且每个 $F_i$ 在 $X$ 中的\Term{closure}{闭包} $\overline F_i$ 都是 $X$ 的\Term{analytic-subspace}{解析子空间}。
\SourcePage{137}

(b) 各\Term{open-stratum}{开层片} $F_i-F_{i+1}$ 在 $D^*$ 上\Term{smooth-morphism}{光滑}。

(c) $\Sigma$ 满足 Whitney 条件 (a) 和 (b)。

对于这样的\Term{stratification}{分层} $\Sigma$，各 $F_i-F_{i+1}$ 上存在\Term{integrable-connection}{可积联络} $\nabla_i$，使得：

(i) 在 $D^*$ 上\Term{local-scheme}{局部}地，$\nabla_i$ 定义 $F_i-F_{i+1}$ 的一个\Term{trivialization}{平凡化}；

(ii) 这些\Term{trivialization}{平凡化}可以\Term{gluing}{粘合}，给出 $X^*/D^*$ 的\Term{local-scheme}{局部}\Term{trivialization}{平凡化}。

最后，若 $(Y_\alpha)$ 是 $X$ 的\Term{finite}{有限}族闭\Term{analytic-subspace}{解析子空间}，则进一步缩小 $D$ 后，可以选取 $\Sigma$，使每个 $Y_\alpha^*$ 都是\Term{open-stratum}{开层片} $(F_i-F_{i+1})$ 的若干\Term{connected-component}{连通分支}之并。此时，上述\Term{connection}{联络} $\nabla$ 给出 $(X,(Y_\alpha))$ 的\Term{local-scheme}{局部}\Term{trivialization}{平凡化}。
\end{SGAresult}

\textbf{1.3.6.} 当 $X$ 在 $D$ 上\Term{proper-morphism}{固有}且下述条件成立时，(1.3.3.2) 对以\Term{sheaf}{层}\Term{complex}{复形} $K$ 为系数的\Term{hypercohomology}{超上同调}也成立：

$(*)$ 对每个 $i\in\mathbf Z$，$D^*$ 可由\Term{open-subset}{开集} $U$ \Term{covering}{覆盖}，使得存在与到 $U$ 的\Term{projection}{投影}相容的\Term{isomorphism}{同构} $\alpha:f^{-1}(U)\xrightarrow{\sim}Y\times U$，以及 $Y$ 上的\Term{sheaf}{层} $G$，使 $\mathcal H^i(K)$ 在 $f^{-1}(U)$ 上\Term{isomorphism}{同构}于 $\alpha^*\operatorname{pr}_1^*G$。

在此假设下，1.2.1 的三角给出\Term{long-exact-sequence}{长正合列}
\[
 \cdots\longrightarrow\mathbb H^i(X_0,K)
 \longrightarrow\mathbb H^i(X_t,K)
 \longrightarrow\mathbb H^i(X_0,R\Phi_t(K))
 \longrightarrow\cdots.\tag{1.3.6.1}
\]
由 Thom \Term{theorem}{定理}，若 $X$ 存在一个复解析\Term{stratification}{分层}，使各 $\mathcal H^i(K)$ 在每个\Term{open-stratum}{开层片}上都是\Term{locally-constant-sheaf}{局部常值层}，则缩小 $D$ 后，条件 $(*)$ 成立。

\textbf{1.4. 一个特殊情形。}

本小节不会在本讨论班后续部分使用。这里把第 XIII 讲引言中略述的一个构造
\SourcePage{138}
\Term{formalization}{形式化}。

\textbf{1.4.1.} 设 $D$ 是一个已\Term{orientation}{定向}的\Term{disc}{圆盘}，$0\in D$，$D^*=D-\{0\}$。$\mathscr D^*$ 上的一个\Term{topological-space}{拓扑空间} $\mathscr X^*$ 是一个\Term{functor}{函子}
\[
 (D^*\text{ 的万有覆盖})\longrightarrow(\text{拓扑空间}).
\]
一旦选定一个\Term{universal-covering}{万有覆盖} $\widetilde D^*$，这就是一个带有 $\pi_1(D^*)$ \Term{action}{作用}的\Term{topological-space}{拓扑空间} $X_\eta$；当 $\widetilde D^*$ 是 $(D^*,t)$ 的\Term{universal-covering}{万有覆盖}时，把它记为 $\mathscr X_t^*$。$\mathscr X^*$ 上的一个\Term{sheaf}{层}是一个\Term{functor}{函子}
\[
 (D^*\text{ 的万有覆盖})\longrightarrow(\text{带有一个层的拓扑空间}),
\]
也就是 $X_\eta$ 上的一个 $\pi_1(D^*)$-\Term{equivariant-sheaf}{等变层}。

给 $\mathscr X^*$ 配上 $D^*$ 上的\Term{topological-space}{拓扑空间}
\[
 \beta^*\mathscr X^*=X_\eta\times\widetilde D^*/\pi_1(D^*).
\]
$\mathscr X^*$ 上的每个\Term{sheaf}{层} $F$ 通过取商，在 $\beta^*\mathscr X^*$ 上定义一个\Term{sheaf}{层} $\beta^*F$，由此得到\Term{topos}{拓扑斯}的\Term{commutative-diagram}{交换图}
\[
\begin{tikzcd}
\beta^*\mathscr X^*\arrow[r]\arrow[d]&\mathscr X^*\arrow[d]\\
D^*\arrow[r,"\beta"]&\mathscr D^*.
\end{tikzcd}
\]

\textbf{1.4.2.} $\mathscr D$ 上的一个\Term{topological-space}{拓扑空间} $\mathscr X$ 由下列资料组成：

(a) $\mathscr D^*$ 上的一个\Term{topological-space}{拓扑空间} $\mathscr X^*$；

(b) 一个\Term{topological-space}{拓扑空间} $X_0$；

(c) 一个 $\pi_1(D^*)$ \Term{invariant}{不变}的\Term{morphism}{态射} $\mathrm{sp}:X_\eta\to X_0$（$\pi_1(D^*)$ 在 $X_0$ 上的\Term{action}{作用}是平凡的）。

$\mathscr X$ 上的一个\Term{sheaf}{层}由 $X_0$ 上的一个\Term{sheaf}{层} $F_0$、$\mathscr X^*$ 上的一个\Term{sheaf}{层}（由 $F_\eta$ 给出），以及一个\Term{equivariant}{等变}\Term{morphism}{态射}
\[
 \mathrm{sp}:\mathrm{sp}^*F_0\longrightarrow F_\eta
 \qquad(\text{或 }F_0\longrightarrow\mathrm{sp}_*F_\eta)
\]
组成。有\Term{commutative-diagram}{交换图}
\SourcePage{139}
\[
\begin{tikzcd}
X_0\arrow[r,hook]\arrow[d]&\mathscr X\arrow[d]&\mathscr X^*\arrow[l,hook']\arrow[d]\\
\{0\}\arrow[r,hook]&\mathscr D&\mathscr D^*\arrow[l,hook'].
\end{tikzcd}
\]
此外，将前述构造的\Term{functoriality}{函子性}应用于
\[
 (X_\eta,X_0,\mathrm{sp})\longrightarrow(X_0,X_0,\mathrm{Id}),
\]
得到一个\Term{topos}{拓扑斯}\Term{morphism}{态射} $\Psi:\mathscr X\to X_0\times\mathscr D$。有
\[
 \Psi(F_0,F_\eta,\mathrm{sp})=(F_0,\mathrm{sp}_*F_\eta,\mathrm{sp}).
\]
此\Term{morphism}{态射}诱导 $\Psi_\eta:\mathscr X^*\to X_0\times\mathscr D^*$。

\textbf{1.4.3.} 对 $\mathscr D$ 上的\Term{topological-space}{拓扑空间} $\mathscr X$，按如下方式配上 $D$ 上的\Term{topological-space}{拓扑空间} $X=\beta^*\mathscr X$。在 $X^*=\beta^*\mathscr X^*$ 中添入位于 $0$ 之上的\Term{fibre}{纤维} $X_0$；对于 $x_0\in X_0$，取下列\Term{set}{集合}作为一个\Term{neighbourhood-basis}{基本邻域系}：
\[
 V\cup\bigl(\mathrm{sp}^{-1}(V)\times\widetilde U^*/\pi_1(D^*)\bigr),
\]
其中 $V$ 是 $X_0$ 中 $x_0$ 的\Term{neighbourhood}{邻域}，$U$ 是 $D$ 中 $0$ 的\Term{neighbourhood}{邻域}，$\widetilde U^*=U\times_D\widetilde D^*$。以 $j$（相应地 $i$）表示 $X^*$（相应地 $X_0$）到 $X$ 的包含。对 $\mathscr X^*$ 上的一个\Term{sheaf}{层} $F$，有
\[
 i^*j_*\beta^*F=(\mathrm{sp}_*F_\eta)^{\pi_1(D^*)}.
\]
从而得到\Term{topos}{拓扑斯}的\Term{commutative-diagram}{交换图}
\[
\begin{tikzcd}
X\arrow[r,"\beta"]\arrow[d]&\mathscr X\arrow[d]\\
D\arrow[r,"\beta"]&\mathscr D,
\end{tikzcd}
\]
使得对 $\mathscr X$ 上的每个\Term{sheaf}{层} $F$，都有
\[
 \Psi(F)\xrightarrow{\sim}\Psi(\beta^*F).\tag{1.4.3.1}
\]

\begin{SGAresult}{命题}{1.4.4}
\Term{arrow}{箭头} (1.4.3.1) \Term{derive}{导出}为一个\Term{isomorphism}{同构}
\[
 R\Psi\xrightarrow{\sim}R\Psi\circ\beta^*,
\]
\SourcePage{140}
其两端为从 $\mathrm D^+(\mathscr X)$ 到 $\mathrm D^+(X_0\times\mathscr D)$ 的\Term{functor}{函子}。
\end{SGAresult}
\begin{proof}
设 $\overline j$ 为
\[
 \widetilde X^*=X\times_D\widetilde D^*\simeq X_\eta\times\widetilde D^*
 \quad\text{到}\quad
 \overline X=X\times_D\widetilde D
\]
的包含。此命题归结为验证：对 $X_\eta$ 上的每个\Term{flasque-sheaf}{松弛阿贝尔层} $F$，$\operatorname{pr}_1^*F$ 对\Term{functor}{函子} $\overline j_*$ 是\Term{acyclic}{无上同调}的。这个断言又归结为以下熟知的事实（取 $B=\widetilde D^*$），它可用一个\Term{homotopy}{同伦}论证证明。

设 $F$ 是\Term{topological-space}{空间} $Y$ 上的一个阿贝尔\Term{sheaf}{层}，$B$ 是一个\Term{disc}{圆盘}。有
\[
 H^i(Y,F)\xrightarrow{\sim}H^i(Y\times B,\operatorname{pr}_1^*F).
\]
\end{proof}

\textbf{1.4.5.} 命题 1.4.4 给出计算 $X$ 上形如 $\beta^*F$ 的\Term{sheaf}{层}的 $R\Psi$ 的方法，实质上把 $R\Psi$ 的计算化为 $R\mathrm{sp}_*$ 的计算。

\textbf{1.4.6.} 后面我们要考虑如下情形：$X_\eta$ \Term{locally-compact}{局部紧}，而 $X_0$ 由 $X_\eta$ 把一个紧\Term{subset}{子集} $V\subset X_\eta$ 压缩为一个点 $0$ 而得；$\mathrm{sp}$ 是压缩映射，$\pi_1(D^*)$ 在 $V$ 之外的\Term{action}{作用}是平凡的。\Term{proper-morphism}{固有态射} $\mathrm{sp}$ 的\Term{base-change}{换基}\Term{theorem}{定理}于是给出：对 $K\in\operatorname{Ob}\mathrm D^+(\mathscr X^*)$，
\[
 R\Psi_\eta(\beta^*K)_0=R\Gamma(V,K),\tag{1.4.6.1}
\]
\[
 R\Gamma_{\{0\}}R\Psi_\eta(\beta^*K)=R\Gamma_V(X_\eta,K)\tag{1.4.6.2}
\]
（\Term{cohomology-with-support}{带支撑上同调}）。

对 $\mathscr X$ 上的 $K$，若 $\mathrm{sp}:\mathrm{sp}^*K_0\to K_\eta$ 在 $V$ 之外为\Term{quasi-isomorphism}{拟同构}，则 $R\Phi(K)$ 在 $0$ 之外为零。特别地，对 $X$ 上的一个\Term{constant-sheaf}{常值层} $\Lambda$，有
\[
 R^i\Phi(\Lambda)=\widetilde H^i(V,\Lambda),\qquad\text{集中于 }0
\]
\SourcePage{141}
（\Term{reduced-cohomology}{约化上同调群}）。

\textbf{1.4.7.} 我们关心的情形是 $X_\eta$ 为\Term{differentiable-manifold}{微分流形}，$V\subset X_\eta$ 是一个内部为 $V^\circ$ 的\Term{manifold-with-boundary}{带边子流形}；此时 $X-\{0\}$ 是一个\Term{differentiable-manifold}{微分流形}，由 $\beta^*X_\eta$ 与 $D$ 上\Term{fibre}{纤维}为 $X_\eta-V$ 的常\Term{topological-fibration}{纤维丛}\Term{gluing}{粘合}而成。映射
\[
 H_c^i(V^\circ,\Lambda)\xrightarrow{\sim}H_V^i(X_\eta,\Lambda)\tag{1.4.7.1}
\]
是\Term{isomorphism}{同构}。这些群也记为 $H^i(V\bmod\partial V,\Lambda)$。\Term{isomorphism}{同构} (1.4.7.1) 定义一个\Term{isomorphism}{同构}
\[
 R\Gamma_c(V^\circ,\Lambda)\xrightarrow{\sim}
 R\Gamma_{\{0\}}R\Psi_\eta(\Lambda),\tag{1.4.7.2}
\]
而\Term{variation}{变差}由一个以 $V$ 的\Term{reduced-cohomology}{约化上同调}为源的\Term{morphism}{态射}定义（用第 XIII 讲 2.4.5 的方法，并经由 (1.4.7.2)）：
\[
 \operatorname{Var}:R\widetilde\Gamma(V,\Lambda)\longrightarrow
 R\Gamma_c(V^\circ,\Lambda).\tag{1.4.7.3}
\]
$T$ 诱导\Term{space}{空间}对 $(V,\partial V)$ 的一个\Term{automorphism}{自同构}，仍记为 $T$，且 $T$ 在 $\partial V$ 上的\Term{restriction}{限制} $\partial T$ 为恒等。\Term{variation}{变差} (1.4.7.3) 应当与赋予此\Term{automorphism}{自同构} $T$ 的经典\Term{variation}{变差}\Term{morphism}{态射}（3.1.9）一致。3.1.10 将证明这一类型的一个命题。

\SGAsection{2}{比较定理}

\textbf{2.1.} 对每个 $\mathbf C$ 上的\Term{finite-type}{有限型}\Term{scheme}{概形} $X$，我们用 $X_{\mathrm{cl}}$ 表示带有通常\Term{topology}{拓扑}的\Term{topological-space}{空间} $X(\mathbf C)$。

设 $S$ 是 $\mathbf C$ 上的一条\Term{smooth-morphism}{光滑}\Term{curve}{曲线}，即 $\mathbf C$ 上的\Term{smooth-morphism}{光滑}\Term{separated-morphism}{分离}\Term{finite-type}{有限型}、纯一维\Term{scheme}{概形}。设 $0\in S(\mathbf C)$，$f:X\to S$ 是一个\Term{separated-morphism}{分离}\Term{finite-type}{有限型} $S$-\Term{scheme}{概形}，$\Lambda$ 是一个\Term{noetherian-ring}{诺特}\Term{torsion-ring}{挠环}（例如 $\mathbf Z/n\mathbf Z$）。置 $X_0=f^{-1}(0)$。

\SourcePage{142}
\textbf{2.2.} 以 $S^h$ 表示 $S$ 在 $0$ 处的\Term{henselization}{Hensel 化}，它是一个\Term{henselian-trait}{亨泽尔离散赋值环谱}；以 $\eta$ 表示 $S^h$ 的\Term{generic-point}{泛点}。有\Term{functor}{函子}
\[
 \begin{aligned}
 \Psi:(X\text{ 上的层})&\xrightarrow{\text{限制}}(X\times_S S^h\text{ 上的层})\\
 &\xrightarrow{\Psi}(X_0\times S^h\text{ 上的层}).
 \end{aligned}
\]
$\Psi$ 的\Term{derived-functor}{导出函子}是复合
\[
 R\Psi:\mathrm D(X,\Lambda)\longrightarrow\mathrm D(X\times_S S^h,\Lambda)
 \xrightarrow{R\Psi}\mathrm D(X_0\times S^h,\Lambda).
\]

\textbf{2.3.} 设 $D\subset S_{\mathrm{cl}}$ 是包含 $0$ 的一个\Term{disc}{圆盘}。复合\Term{functor}{函子} $\Psi_{\mathrm{cl}}$（或简记为 $\Psi$）
\[
\begin{aligned}
 \Psi_{\mathrm{cl}}:(X_{\mathrm{cl}}\text{ 上的层})
 &\xrightarrow{\text{限制}}(X_{\mathrm{cl}}\times_{S_{\mathrm{cl}}}D\text{ 上的层})\\
 &\xrightarrow{\Psi}((X_0)_{\mathrm{cl}}\times\mathscr D\text{ 上的层})
\end{aligned}
\]
与 $D$ 的选择无关（在唯一\Term{isomorphism}{同构}意义下；回想把 $D$ 换成较小的\Term{disc}{圆盘}时，$\mathscr D$ \Term{invariant}{不变}）。它的\Term{derived-functor}{导出函子}是复合
\[
 R\Psi_{\mathrm{cl}}:\mathrm D(X_{\mathrm{cl}},\Lambda)
 \longrightarrow\mathrm D(f_{\mathrm{cl}}^{-1}(D),\Lambda)
 \xrightarrow{R\Psi}\mathrm D(X_0\times\mathscr D,\Lambda).
\]

\textbf{2.4.} 类似地定义
\[
\begin{aligned}
 R\Psi_\eta&:\mathrm D(X-X_0,\Lambda)\longrightarrow\mathrm D(X_0\times\eta,\Lambda),\\
 R\Psi_{\eta\mathrm{cl}}&:\mathrm D((X-X_0)_{\mathrm{cl}},\Lambda)
 \longrightarrow\mathrm D(X_0\times\mathscr D^*,\Lambda).
\end{aligned}
\]

\textbf{2.5.} 构造 1.1.11 定义一个从 $\mathscr D^*$ 到 $\eta$ 的\Term{topos}{拓扑斯}\Term{morphism}{态射}，它延拓为
\[
 \epsilon:\mathscr D\longrightarrow S^h.
\]
此\Term{morphism}{态射}与 $\epsilon:(X_0)_{\mathrm{cl}}\to X_0$（SGA 4，第 XI 讲第 4 节）一起定义
\[
 \epsilon:(X_0)_{\mathrm{cl}}\times\mathscr D\longrightarrow X_0\times S^h.
\]
\SourcePage{143}
若 $\widetilde D^*$ 是 $D^*$ 的一个\Term{universal-covering}{万有覆盖}，$k(\overline\eta)$ 是相应的 $k(\eta)$ 的\Term{algebraic-closure}{代数闭包}（1.1.10），则有
\[
 \epsilon^*(F_0,F_\eta,\alpha)
 =(\epsilon^*F_0,\epsilon^*F_{\overline\eta},\epsilon^*\alpha).
\]

\textbf{2.6.} 设 $\widetilde D^*$ 是 $D^*$ 的一个\Term{universal-covering}{万有覆盖}，$k(\overline\eta)$ 是相应的 $k(\eta)$ 的\Term{algebraic-closure}{代数闭包}，$\overline S$ 是 $S^h$ 在 $k(\overline\eta)$ 中的\Term{normalization}{正规化}。把 $\overline S$ 写成\Term{affine-scheme}{仿射} $S$-\Term{scheme}{概形} $S_i$ 的\Term{inverse-limit}{投射极限}；对于每个 $i$，把 $D$ 换成较小的\Term{disc}{圆盘} $D_i$ 后，有自然的 $S_{\mathrm{cl}}$-\Term{topological-space}{空间}\Term{morphism}{态射}
\[
 \widetilde D_i\xrightarrow{a_i}(S_i)_{\mathrm{cl}}.
\]
设 $j_i:U_i\hookrightarrow S_i$ 是 $S_i$ 中 $\overline S$ 的\Term{closed-point}{闭点}之像 $0$ 的补\Term{open-subset}{开集}，$i:0\hookrightarrow S_i$。

设 $F$ 是 $X$ 上的\Term{sheaf}{层}，$F_{\mathrm{cl}}=\epsilon^*F$。仍用 $j_i$ 表示 $X\times_S U_i$ 到 $X\times_S S_i$ 的包含，用 $i$ 表示 $X_0$ 到 $X\times_S S_i$ 的包含，并仍用 $F$ 表示 $F$ 在 $X\times_S S_i$ 上的\Term{inverse-image}{逆像}。有
\[
 \Psi(F)=\left(i^*F,\varinjlim_i i^*j_{i*}j_i^*F,
                       \varinjlim_i(\varphi_i)\right).
\]
\Term{functor}{函子} $\epsilon^*$ 与\Term{inverse-image}{逆像}\Term{functor}{函子}交换。因此\Term{morphism}{态射} $\epsilon^*j_{i*}\to(j_{i\mathrm{cl}})_*\epsilon^*$ 给出
\[
 c_i:(i^*j_{i*}j_i^*F)_{\mathrm{cl}}
 \longrightarrow i^*j_{i*}j_i^*(F_{\mathrm{cl}}).
\]
设 $\overline j$（相应地 $\overline j_i$）是 $\widetilde D^*$（相应地 $\widetilde D_i^*$）到 $\widetilde D$（相应地 $\widetilde D_i$）的包含，也用它们表示由此得到的包含
\[
 X\times_S\widetilde D^*\hookrightarrow X\times_S\widetilde D,
 \qquad X\times_S\widetilde D_i^*\hookrightarrow X\times_S\widetilde D_i.
\]
\Term{adjunction}{伴随}\Term{arrow}{箭头} $\mathrm{Id}\to a_{i*}a_i^*$ 定义
\[
 i^*j_{i*}j_i^*F_{\mathrm{cl}}
 \longrightarrow i^*\overline j_{i*}\overline j_i^*F_{\mathrm{cl}}
 =i^*\overline j_*\overline j^*F_{\mathrm{cl}},
\]
故取\Term{limit}{极限}得到
\[
 c_\eta:\Psi_{\overline\eta}(F)_{\mathrm{cl}}
 \longrightarrow\Psi_{\mathrm{cl}}(F_{\mathrm{cl}}).
\]
此\Term{morphism}{态射}定义\Term{comparison-morphism}{比较态射}
\SourcePage{144}
\[
 c:\epsilon^*\Psi(F)\longrightarrow\Psi_{\mathrm{cl}}(\epsilon^*F).\tag{2.6.1}
\]
类似地定义
\[
 c_\eta:\epsilon^*\Psi_\eta(F)\longrightarrow
 \Psi_{\eta\mathrm{cl}}(\epsilon^*F).\tag{2.6.2}
\]

\textbf{2.7.} \Term{morphism}{态射} (2.6.1) 和 (2.6.2) \Term{derive}{导出}为
\[
\begin{aligned}
 c&:\epsilon^*R\Psi\longrightarrow R\Psi_{\mathrm{cl}}\circ\epsilon^*,\\
 c_\eta&:\epsilon^*R\Psi_\eta\longrightarrow R\Psi_{\eta\mathrm{cl}}\circ\epsilon^*.
\end{aligned}
\]

\begin{SGAresult}{\Term{theorem}{定理}}{2.8}
设 $K\in\operatorname{Ob}\mathrm D(X,\Lambda)$（相应地 $\operatorname{Ob}\mathrm D(X-X_0,\Lambda)$）是一个 $\Lambda$-\Term{module}{模}\Term{sheaf}{层}的\Term{complex}{复形}，其\Term{cohomology-sheaf}{上同调层} $\mathcal H^i(K)$ 都是\Term{constructible-sheaf}{可构造的}。则\Term{comparison-morphism}{比较态射}
\[
 c:\epsilon^*R\Psi(K)\longrightarrow R\Psi_{\mathrm{cl}}(\epsilon^*K)
\]
（相应地
\[
 c:\epsilon^*R\Psi_\eta(K)\longrightarrow R\Psi_{\eta\mathrm{cl}}(\epsilon^*K)
\]
）是\Term{isomorphism}{同构}。
\end{SGAresult}
\begin{proof}
只需对 $R\Psi_\eta$ 证明此断言。又因 (2.7) 两端都是\Term{cohomological-dimension}{上同调维数}\Term{finite}{有限}的\Term{triangulated-functor}{三角函子}，只需处理 $K$ 归结为一个\Term{sheaf}{层} $F$（置于零次）的情形。证明与第 XIII 讲\Term{finiteness-theorem}{有限性定理} 2.3.1 的证明完全类似：$R\Psi$ 与 $R\Psi_{\mathrm{cl}}$ 具有相同的\Term{covariant}{协变}、\Term{contravariant}{反变}性质，故由该处的\Term{devissage}{分解步骤}（其中允许缩小 $S$），可归结为 $X$ \Term{regular}{正则}、$f^{-1}(0)$ 是 $X$ 中的一个\Term{normal-crossings-divisor}{正规交叉除子}且 $F=\mathbf Z/\ell$ 的情形。然后把第 I 讲 3.3 与其\Term{transcendental}{超越}对应结果进行比较。
\end{proof}
